\section{Related Work}
\label{sec:related}

\subsection{Automated environment construction from papers}

RSI-Forge introduced a multi-agent pipeline that turns published papers into
executable environments for self-improvement \cite{rsi-forge}. Three agents
coordinate \emph{construction, reproduction, and review}: a construction agent
generates code and task definitions; a reproduction agent independently
reimplements the method to establish a baseline score; and a review agent checks
task validity and evaluator reliability. The pipeline produced 210 environments
across 18 fields, including 90 reviewed by independent human domain experts.

In a key experiment, four models performed three sequential attempts across 120
environments, inheriting code and notes but \emph{without weight updates}. At
least one model improved after the first attempt in \textbf{84\% of
environments}, and in \textbf{68 of 120} environments model performance exceeded
the reproduced baseline---demonstrating that these environments provide genuine
improvement signals. \sys\ adopts the same three roles and adds two: a
\emph{Curator} that screens the literature before any generation cost is
incurred, and a \emph{Packager} that promotes validated environments into the
standardized \harbor\ layout (\S\ref{sec:method}).

\subsection{\harbor\ and containerized agent evaluation}

\harbor\ is a framework from the creators of Terminal-Bench for evaluating and
optimizing agents and language models in container environments \cite{harbor}.
Terminal-Bench~2.0 emphasizes \emph{task validation quality} to ensure reliable
and reproducible agent evaluation \cite{terminalbench}. \harbor\ has been
adopted by external benchmarks; for example, CompileBench migrated to \harbor\
to standardize agent evaluation for compiling open-source projects
\cite{compilebench}.

\subsection{Life-science agent benchmarks}

DrugDiscoveryBench measures how AI agents handle computational tasks in
early-stage drug discovery, with 82 verifiable, domain-expert-curated tasks
spanning target identification to structure--activity analysis
\cite{drugdiscoverybench}. BixBench provides over 50 real-world scenarios and
nearly 300 open-ended questions in computational biology, with model accuracies
as low as 17\% in complex bioinformatics workflows \cite{bixbench}.
ScienceAgentBench includes 102 scientific data-analysis tasks across
computational biology, geoscience, psychology, and materials science
\cite{scienceagentbench}. LAB-Bench provides a further set of literature-scale
biological reasoning tasks \cite{labbench}. All of these are hand-built; none is
generated from the literature automatically.

\subsection{Reinforcement learning with verifiable and rubric-based rewards}

\textbf{RLVR (Reinforcement Learning with Verifiable Rewards)} has emerged as a
paradigm for equipping LLMs with agentic capabilities through verifiable reward
signals \cite{rlvr}. Recent work extends RLVR through group verification-based
policy optimization (GVPO) \cite{gvpo} and entropy-aware training
\cite{entropyaware}. \textbf{Rubrics as Rewards (RaR)} addresses the limitation
of RLVR in domains without a single standard answer by using structured,
checklist-style rubrics as interpretable reward signals \cite{rar}. RaR achieves
relative improvements of up to 31\% on HealthBench and 7\% on GPQA-Diamond over
LLM-as-judge baselines. Stage~4 of our protocol (\S\ref{sec:stage4}) is designed
to use both regimes, because life-science tasks split cleanly into those with a
verifiable metric (AUROC, F1, RMSE, peak-calling accuracy) and those requiring a
rubric.

\subsection{Held-out benchmarks and contamination}

Benchmark contamination---where test data appears in training corpora---can
inflate performance estimates by as much as 10 percentage points
\cite{contamination}. Truly held-out evaluation requires data the model has
\emph{definitively never encountered}, often approximated through time-based
splits or newly released examination questions. Our Stage~4 therefore trains only
on environments derived from papers published before 2024 and evaluates on
environments derived from 2025--2026 papers, \emph{locked before training
begins} and audited for n-gram overlap (\S\ref{sec:stage4}, claim~\cid{P10}).
